\section{Chat Proving Ground y API Contract}

La última sección conecta la investigación con el producto. En la Chat experience, el proveedor controla la UI, el system prompt, las tool, el rollback y la user education, por lo que puede experimentar con rapidez; en cambio, una vez que un cliente incorpora la API a production, schema, model behavior, latency, error, pricing y retirement se convierten en un contract externo. Como se dijo en clase, «APIs are forever» no significa literalmente que nunca cambie, sino que el costo de cada cambio recae en una gran cantidad de sistemas de clientes.

\subsection{Del experimento a la versioned primitive}

Después de validar en Chat la PDF upload, el long context o el tool behavior, el equipo necesita convertir las restricciones implícitas de UX en API schema, auth, quota, idempotency, error model, version y docs. La \term{API deprecation} (retiro de la recomendación de uso) consiste en anunciar que un modelo o un parámetro ya no se recomienda, y en ofrecer replacement, migration guidance y retirement date; después del retirement, las solicitudes fallan, por lo que los clientes deben probar y migrar dentro de ese plazo.

\lecturefigure{16-chat-to-api.png}{Chat se encarga de la experimentación rápida; la API publica las capacidades maduras como un contract compatible a largo plazo.}{Subtítulos locales 39:00--41:56; Anthropic API lifecycle docs; redibujo conceptual.}

\begin{knowledgebox}{Lectura de la figura: qué evidencia debe completarse antes del release}
El Chat experiment observa primero la usefulness, la seguridad y las failure; la evidence da lugar a los acceptance criteria; el API release fija schema/version/SLO; el lifecycle ofrece los estados active, legacy, deprecated y retired, junto con la migration. La continuity del cliente puede importar más que el hecho de que «el modelo más reciente sea más inteligente»; por ello, el replacement debe someterse a pruebas de regresión con el workload real, y no basta con fijarse en el benchmark.
\end{knowledgebox}

\teachervoice{Mann pone como ejemplo que modelos antiguos de Claude siguen en algún entorno de producción para ilustrar la API inertia: es posible que el cliente no tenga recursos de ingeniería para migrar, o que haya incorporado el comportamiento anterior en sus procesos de cumplimiento normativo. La clase advierte a los frontier lab que publicar una API equivale a aceptar responsabilidades de support, compatibility y deprecation.}

\begin{warningbox}{Actualizar el modelo no es un binary upgrade ordinario}
Un modelo nuevo puede cambiar refusal, format, tool call, latency, token usage y la nondeterministic distribution. La migration requiere golden tasks, safety regression, cost/SLO test, shadow traffic, rollback y customer communication; cambiar en silencio un alias con el mismo nombre hace que sea difícil atribuir la causa de un incidente.
\end{warningbox}

\subsection{Resumen de la sección}

Es razonable que Chat y la API avancen a velocidades distintas. El proving ground ayuda a descubrir comportamientos; el API contract exige versionado, capacidad de prueba y capacidad de migración. La disciplina de producto es el último paso para que una frontier capability se convierta en infraestructura.
